\section{Mercy and the purpose of mercy}
Chrysostom and Augustine expose different failures in people whose lives can look commendable. Chrysostom asks whether an admired discipline has left the neighbor without help. Augustine asks whether help already given proceeds from charity. A person can therefore need the first examination, the second, or both. Their compatibility as moral demands does not make their explanations of the lamps and sellers identical.

\begin{comparisontable}{Question}{Mercy practiced}{Charity within}{Consequence}
What is missing? & Chrysostom: humanity and assistance despite a demanding virtue. & Augustine: inward charity despite observable good works. & Examine both the help withheld and the purpose of help given.\\
Who sells oil? & The poor, who should have been approached during life. & Flatterers, whose praise has sustained an outward reputation. & The symbolic identifications differ; they are not one patristic consensus.\\
Why can oil not be shared? & A history of merciful action cannot be transferred after death. & Another person's inward testimony cannot be supplied at judgment. & The final refusal gives no rule for withholding present help.\\
How is Christ encountered? & Gregory's hospitality explains recognition through acting on the word. & The same hospitality tests an inward claim through an actual deed. & The Luke Communion alternative concerns welcomed and recognized presence.\\
Why seek help after Communion? & Mercy requires continued availability and perseverance. & The purpose of good action requires continued purification. & Both interpretations depend on the Spirit's sustaining gift.\\
\end{comparisontable}

The two accounts also agree where the story itself demands agreement. All sleep, so ordinary bodily weariness does not distinguish the foolish. The cry arrives without a usable schedule, so a plan to become ready later remains uncertain. The door closes, so preparation has a consequence beyond the moment in which it is undertaken. Wisdom's initiative and Christ's resurrection give the response its beginning and its hope. Within that hope, Chrysostom directs attention outward to a person awaiting mercy, and Augustine directs attention to the charity with which that mercy is offered.
